\documentclass{article}
\usepackage{amsmath, amssymb, amsthm}
\usepackage{hyperref}
\usepackage{geometry}
\geometry{margin=1in}

\title{Erd\H{o}s Problem \#640}
\date{Last updated: 2025-08-31}
\author{Source: erdosproblems.com}

\begin{document}
\maketitle

\noindent\textbf{Status:} OPEN \quad \textbf{No prize}

\noindent\textbf{Tags:} graph theory, chromatic number

\medskip
\noindent\textbf{Problem Statement:}

Let $k\geq 3$. Does there exist some $f(k)$ such that if a graph $G$ has chromatic number $\geq f(k)$ then $G$ must contain some odd cycle whose vertices span a graph of chromatic number $\geq k$?

\bigskip
\noindent\textbf{Remarks:}

A problem of Erd\H{o}s and Hajnal. 

This is trivial when $k=3$, since any non-bipartite graph contains an odd cycle, and all odd cycles have chromatic number $3$.

Steiner has observed in the comments that this is equivalent to the problem in which we replace 'odd cycle' with 'path'.

\bigskip
\noindent\small{Source: \url{https://www.erdosproblems.com/640}}
\end{document}
